% TOP_PROBLEM: {"id": 3248, "title": "Bounding the on-line choice number in terms of the choice number", "category_id": 3, "category": {"id": 3, "name": "graph_theory", "display_name": "Graph Theory", "description": "Problems involving graphs, networks, and their properties.", "slug": "graph-theory", "order_index": 3, "created_at": "2026-07-31T15:26:25.671Z"}, "status": "open", "problem_number": "OPG-48583", "set_id": 12}
% FIRST_POSTED: 2026-10-01
% ENTRY_KIND: statement_only
% UnsolvedMath ID 3248; source catalogue code OPG-48583
% !TeX program = lualatex
% Definition and sources only; this file is not an LLM solution attempt.
\documentclass[11pt,a4paper]{article}
\usepackage[margin=25mm]{geometry}
\usepackage{fontspec}
\setmainfont{lmroman10-regular.otf}[BoldFont=lmroman10-bold.otf,
 ItalicFont=lmroman10-italic.otf,BoldItalicFont=lmroman10-bolditalic.otf]
\usepackage{amsmath,amssymb,mathtools}
\usepackage{unicode-math}
\setmathfont{latinmodern-math.otf}
\AtBeginDocument{\let\setminus\smallsetminus}
\usepackage{xurl,hyperref}
\hypersetup{colorlinks=true,urlcolor=blue}
\setcounter{secnumdepth}{0}
\setlength{\parindent}{0pt}
\setlength{\parskip}{5pt}
\setlength{\emergencystretch}{3em}
\begin{document}
\section*{Bounding the on-line choice number in terms of the choice number}\label{3248}
{\small UnsolvedMath ID 3248 · OPG-48583 · Graph Theory · OpenGarden}
\subsection{Definitions and mathematical statement}
For a finite simple graph $G$, its choice number $\operatorname{ch}(G)$ is
the least $k$ such that every assignment of $k$ permitted colours per
vertex admits a proper colouring from those lists.

Define the online version through the following finite game. Initially
every vertex is uncoloured and has $k$ tokens. In each round Lister marks
a nonempty subset $M$ of the uncoloured vertices and removes one token
from each marked vertex. Painter then selects an independent subset
$I\subseteq M$ and permanently colours those vertices with the fresh
colour assigned to that round. If any vertex remains uncoloured with
zero tokens after Painter's response, Lister wins. Painter wins if all
vertices become coloured. Strategies may depend on the complete history;
neither player sees the other's future decisions.

The online choice number $\operatorname{ch}^{\mathrm{OL}}(G)$ is the least
$k$ for which Painter has a winning strategy against every play by Lister.
The independent-set requirement guarantees a proper colouring because
different rounds use different colours. This is also called the paint
number.

The source question is whether the additive gap is unbounded:
\[
 \forall b\in\mathbb N\ \exists\text{ a finite simple graph }G:
 \quad\operatorname{ch}^{\mathrm{OL}}(G)-\operatorname{ch}(G)\ge b.
\]
The graph may depend on $b$. A negative answer would give a single
absolute additive bound valid for all graphs. This question asks about
the difference, not whether the online parameter is bounded by some
arbitrary function of the offline parameter.
\subsection{Short English statement}
Can the online list-colouring requirement exceed the ordinary choice number by an arbitrarily large additive amount?
\subsection{Sources}
\begin{itemize}
\item[S1] ULAM.AI and the UnsolvedMath Contributors, supplied problems.json, record 3248 (OPG-48583), OpenGarden collection. Catalogue identity and source transcription; CC BY 4.0.
\url{https://huggingface.co/datasets/ulamai/UnsolvedMath}.
\item[S2] Open Problem Garden, Bounding the on-line choice number in terms of the choice number. Original problem and discussion; the mathematical formulation below is expanded from the supplied source record.
\url{http://www.openproblemgarden.org/op/bounding_the_on_line_choice_number_in_terms_of_the_choice_number}.
\end{itemize}
\end{document}
